\section{Introduction}
\label{sec:intro-body}

Random walks on groups are among the most studied objects in probability theory
and harmonic analysis. Given a finitely generated group $G$ with a symmetric
generating set $S$, the simple random walk has transition kernel
\begin{equation}
\label{eq:kernel}
P(x, y) = \frac{1}{|S|} \, \mathbf{1}_{y \in x S},
\end{equation}
and its $n$-step return probabilities $p^{(n)}(e)$ govern mixing, entropy, and
spectral phenomena. A central object is the \emph{spectral radius}
\begin{equation}
\label{eq:rho}
\rho(G, S) = \limsup_{n \to \infty} p^{(n)}(e)^{1/n}.
\end{equation}

For free groups, Kesten proved $\rho = \frac{2\sqrt{|S|-1}}{|S|}$, and it is
known that $\rho = 1$ characterizes amenability \cite{Kesten1959, Woess2000}.
Our aim is a quantitative refinement of these facts for \emph{finite} groups,
where the analogue of $\rho$ is the second-largest eigenvalue $\lambda_1$ of
$P$, and the \emph{spectral gap} is $1 - \lambda_1$.

We work throughout with the normalized counting measure, and we write
$|A|$ for the cardinality of a set $A$. The Cayley graph $\Gamma(G, S)$ is
regular of degree $d = |S|$; its adjacency operator acts on $\ell^2(G)$.
The reader is referred to \cite{Diaconis} for background on random walks on
groups and to \cite{Lubotzky} for expander constructions.

\begin{remark}[Notation]
\label{rem:notation}
All logarithms are natural, and $e$ denotes the identity element of the group
as well as Euler's number; context disambiguates the two.\footnote{This
deliberately ambiguous convention is standard in the literature on random
walks on groups, cf.\ \cite{Woess2000}.}
\end{remark}
